\documentclass[11pt,fleqn]{report}
\usepackage{graphicx}
\usepackage{amssymb}
\usepackage{longtable}
\usepackage[T1]{fontenc}

\usepackage{cml}
\lstnewenvironment{codeXML}
    {\lstset{basicstyle=\small\ttfamily,
             language=XML,
             xleftmargin=0.0in,
             xrightmargin=0.0in,
             backgroundcolor=\color{Orchid!60}, % 60% orchid, 40% white
             escapeinside={\%*}{*)}, %  .tex file characters inside %*...*) are
                                     %  treated as Tex commands.
            }
    }
    { }

\newcommand{\ModelName}{Sim Speed Monitor}
\newcommand{\ModelAuthor}{Gary Turner}
\newcommand{\ModelAffiliation}{Odyssey Space Research}
\newcommand{\ModelDate}{Sept 2022}
\pagestyle{plain}
\usepackage[parfill]{parskip}
\usepackage[margin=1.5cm]{caption}
\usepackage{subfig}

\begin{document}
\titlePage
\clearpage % Reset page number
\pagenumbering{roman}
\tableofcontents % Print table of contents
\vfill
{\Large \textbf {Revision History}}
{\large
\begin{center}
\begin{tabular}{|l||l|l|l|}
\hline
 \textbf{Version} & \textbf{Date} & \textbf{Author} & \textbf{Purpose} \\ \hline
 \hline
1 & Sept 2022 & Gary Turner & Initial Version \\ \hline
\end{tabular}
\end{center}
}

\chapter{Introduction} % Make a "chapter" heading
\pagenumbering{arabic} % Start text with arabic 1
The Sim Speed Monitor model computes the sim-speed at regular intervals.
This is useful for debugging issues identified as sim slowdowns because it
makes it easy to identify the frequency and location within a cycle of where
significant slowdowns occur.

\chapter {Requirements}
The model shall provide the sim-speed at a user-specified frequency.

\chapter {Model Specification}
I start this section by defining two terms:
\begin{itemize}
\item Simulation-time is the clock used by the simulation-engine.

  Aside -- because this model relies on scheduling by the simulation engine,
  it must depend on simulation-time rather than dynamic-time.
  For the vast majority of simulations, simulation-time and dynamic-time are
  equivalent.
\item True-time is the time experienced by the user in the real world; this is
  also known as wall-time, or wall-clock-time.
\item Simulation-speed is the ratio of simulation-time to true-time, it
  indicates how many times faster than real-time the simulation is processing.
  For the Trick simulation engine, we get an average simulation-speed at the
  completion of the run. This model provides the user with outputs for
  monitoring the simulation-speed as it varies during the course of the
  simulation.
\end{itemize}
The model computes the simulation speed as a multiple of wall-clock-time.
The elapsed clock-time between calls to the model's \textit{calculate\_rate}
method can also be accessed.

The model contains a single class, \textit{SimSpeedMonitor}.

This class provides two methods of interest:
\begin{enumerate}
\item \textit{calculate\_rate( double simulation\_dt)} is the main method. The
  one argument to this method is the period of simulation-time at which the
  method is called. This method populates the class member
  \textit{\textbf{sim\_speed}}.
\item \textit{get\_dt()} provides a getter to return the true-time elapsed
  during the most recent computation of \textit{sim\_speed}.
\end{enumerate}


\chapter {User's Guide}
\section {Implementation}
Somewhere in the simulation, instantiate an instance of
\textit{SimSpeedMonitor}. Typically, it would be desirable to monitor several
different frequencies at once so this may be an array of instances:

\begin{codeC}
SimSpeedMonitor  sim_speed[3];
\end{codeC}

\section{Initialization}
The model is self-initializing on the first call to its
\textit{calculate\_rate(...)} method.

\section {Routine Execution}
At some regular interval, call each of these instance's
\textit{calculate\_rate(...)} methods. These could be called at distinct rates:

\begin{codeC}
(0.1, "environment") sim_speed[0].calculate_rate(0.1);
(0.2, "environment") sim_speed[1].calculate_rate(0.2);
(0.5, "environment") sim_speed[2].calculate_rate(0.5);
\end{codeC}

or at distinct points within a cycle:
\begin{codeC}
(0.1, "environment") sim_speed[0].calculate_rate(0.1);
(0.1, "scheduled")   sim_speed[1].calculate_rate(0.1);
(0.1, "effector")    sim_speed[2].calculate_rate(0.1);
\end{codeC}

Important notes:
\begin{enumerate}
\item It is important to note that the value in the argument to
      \textit{calculate\_rate(...)} is identical to the calling period.
\item If trying to identify the period at which the slowdown occurs, having the
      instances run at different rates is most effective.
\item If trying to identify where the slowdown occurs within a cycle, having
      the instances run at different locations within a cycle is most effective.
      In this case, these jobs should run at double the frequency of interest;
      running at the frequency of interest will not return any useful
      information.
\end{enumerate}

\section {Optional Execution}
The regular execution will populate the \textit{sim\_speed} variable. For some
analyses, it may be desirable to also monitor the true-time elapsed between
subsequent calls to \textit{calculate\_rate(...)}. The \textit{get\_dt()}
method can be used to log this value:

\begin{codeC}
double dt = sim_speed[0].get_dt();
\end{codeC}

\section {Configuration}
The model has little available for configuration. As with most CML models, the
instance must be subscribed in order to function.
\begin{codePython}
test.sim_speed[0].subscribe()
\end{codePython}

\chapter{Verification}
The model includes a simple verification simulation, comprising 4 instances of
SimSpeedMonitor, and 2 methods to slow down the simulation, one running at 5Hz
and one at 1Hz. Both slow-down methods execute as ``scheduled'' class jobs.

\begin{center}
\begin{tabular}{|c|c|c|}
\hline
 \textbf{Instance index} & \textbf{Calling Period / s} &
 \textbf{Calling Phase} \\ \hline
 \hline
0 & 0.1 & environment (early) \\ \hline
1 & 0.1 & effector (late)     \\ \hline
2 & 0.5 & environment (early) \\ \hline
3 & 1.0 & environment (early) \\ \hline
\end{tabular}
\end{center}

The following table contains values characteristic of
the simulation-speed measured by the 4 instances.

Note that the actual values from the
verification simulation cannot be committed to the repository for comparison
during regression testing because these values are a function of how quickly
the simulation executes. This is a function of the performance of the machine
which will vary from run to run and machine to machine.

\begin{center}
\begin{tabular}{|c|c|c|c|c|}
\hline
 \textbf{Simulation Time / s} & \textbf{SimSpeed [0]}  & \textbf{SimSpeed [1]} &
 \textbf{SimSpeed [2]} & \textbf{SimSpeed [3]} \\
     & 0.1 s early & 0.1 s late & 0.5 s early & 1.0 s early \\  \hline \hline
0.0  &     0   &     0  &  0  &   0   \\ \hline
0.1  &     3   &  4000  &  0  &   0   \\ \hline
0.2  & 10000   &    10  &  0  &   0   \\ \hline
0.3  &    10   & 10000  &  0  &   0   \\ \hline
0.4  & 10000   &    10  &  0  &   0   \\ \hline
0.5  &    10   & 10000  &  3  &   0   \\ \hline
0.6  & 10000   &    10  &  3  &   0   \\ \hline
0.7  &    10   & 10000  &  3  &   0   \\ \hline
0.8  & 10000   &    10  &  3  &   0   \\ \hline
0.9  &    10   & 10000  &  3  &   0   \\ \hline
1.0  &  3000   &     3  & 10  &  10   \\ \hline
1.1  &     3   & 10000  & 10  &  10   \\ \hline
1.2  & 10000   &    10  & 10  &  10   \\ \hline
1.3  &    10   & 10000  & 10  &  10   \\ \hline
1.4  & 10000   &    10  & 10  &  10   \\ \hline
1.5  &    10   & 10000  &  3  &  10   \\ \hline
1.6  & 10000   &    10  &  3  &  10   \\ \hline
1.7  &    10   & 10000  &  3  &  10   \\ \hline
1.8  & 10000   &    10  &  3  &  10   \\ \hline
1.9  &    10   & 10000  &  3  &  10   \\ \hline
2.0  &  4000   &     3  & 10  &  10   \\ \hline
2.1  &     3   & 10000  & 10  &  10   \\ \hline
2.2  & 10000   &    10  & 10  &  10   \\ \hline
2.3  &    10   & 10000  & 10  &  10   \\ \hline
2.4  & 10000   &    10  & 10  &  10   \\ \hline
2.5  &    10   & 10000  &  3  &  10   \\ \hline
2.6  & 10000   &    10  &  3  &  10   \\ \hline
2.7  &    10   & 10000  &  3  &  10   \\ \hline
2.8  & 10000   &    10  &  3  &  10   \\ \hline
2.9  &    10   & 10000  &  3  &  10   \\ \hline
3.0  &  4000   &     3  & 10  &  10   \\ \hline
\end{tabular}
\end{center}


To identify the calling frequency of the slow-down jobs, we use the instances to
cover a distribution of frequencies.

The instance at index 3, executing at $1 Hz$, shows a
constant speed throughout the simulation. Whatever is causing our slowdown is
happening equally on every $1 Hz$ frame, so must be executing at no less than
$1 Hz$.

The instance at index 2, executing at $2 Hz$ shows a distinct alternating
pattern.  For time \newline $t \in \{[0.5, 1.0), [1.5, 2.0), etc\}$, the
simulation runs relatively slowly when compared to time\newline
$t \in \{[1.0, 1.5), [2.0, 2.5), etc.\}$.
From this information, we can see that some job is having a significant impact
to the simulation speed sometime between the sim-speed detection at $t=1.0$
(when it calculates a high speed), and $t=1.5$, and that this pattern repeats
every second.

The instance at index 0, executing at $10 Hz$ also shows a distinct alternating
pattern, this time at $5 Hz$. At some time between the sim-speed detection at
$t=0.2$
(when it calculates a high speed), and $t=0.3$ (when it calculates a low speed),
we can deduce that a job is executed that produces a significant impact
on average sim speed.

Note also that this higher-frequency job cannot be solely responsible for the
signature seen at $1 Hz$; if the sim-speeds at $2 Hz$ were in ratios of $2:3$,
it could be argued that a 5Hz job was hitting one $2 Hz$ interval twice,
and the next one 3 times. But the ratios are inconsistent with this
inference.

Thus, we can conclude that there are two significant jobs processing within
this simulation, one operating at $5 Hz$, and one at $1 Hz$.  Actually, we
would need to run this for more than 3 seconds to identify the precise
frequency, but this exercise is primarily for illustrative purposes.

To identify the phase of the $5 Hz$ job, we use the two $10 Hz$ speed
calculations
executing at the same timestamps but with different phasing within each cycle.
We see that the instance at index 0, running early in the cycle, computes a
high speed at $t=0.2$, while the instance with index 1, running late in the
cycle, computes a low speed for $t=0.2$. From this information we deduce that
the $5 Hz$ job is executing somewhere in the cycle at $t=0.2$, sandwiched
between the phases at which these two calculations occur -- the first
calculation did not see the impact of this job, but the second calculation
did.

A similar analysis can be run to find where the $1 Hz$ slow-down job executes.

In a large simulation, these jobs can be inserted as desired to systematically
identify the precise top-level function producing the slowdown without having
to resort to a full performance analysis (which can be very
time-consuming).

\end{document}
